%
% teil5.tex -- Anwendung: Die Kepler-Gleichung
%
% (c) 2026 Gian Cavegn, OST Ostschweizer Fachhochschule
%
% !TEX root = ../../buch.tex
% !TEX encoding = UTF-8
%
\section{Anwendung: Die Kepler-Gleichung
\label{bessel:section:kepler}}
\kopfrechts{Kepler-Gleichung}
% Funktion dieses Abschnitts:
% Zweite Anwendung. Hier kommt der historische Ursprung der Bessel-
% Funktionen zur Sprache und zugleich der Algebra-Analysis-Bogen
% des ganzen Buches: Kepler-Gesetze (Geometrie/Algebra der Bahn)
% führen auf eine transzendente Gleichung, deren Lösung eine
% analytische Reihenentwicklung benötigt -- und die Koeffizienten
% sind genau die Bessel-Funktionen.
%
% Wichtig: dieser Abschnitt rechtfertigt auch den Namen "Bessel-
% Funktionen" -- Bessel war Astronom und hat die Funktionen genau
% in diesem Kontext eingeführt.

Die Bessel-Funktionen erhielten ihren Namen nicht durch ihre Rolle in der Wellenoptik, sondern durch eine ältere Anwendung in der Himmelsmechanik.
Friedrich Wilhelm Bessel führte sie 1824 im Zusammenhang mit der Berechnung von Planetenbahnen ein.
% [TODO: Quellenangabe Bessels Originalarbeit von 1824]
Dieser Abschnitt zeigt, wie aus den Kepler'schen Gesetzen eine transzendente Gleichung entsteht, deren Lösung eine Fourier-Reihe mit Bessel-Funktionen als Koeffizienten besitzt.

\subsection{Die Kepler'schen Gesetze
\label{bessel:subsection:keplergesetze}}
% Inhalt:
% Knappe Wiederholung der drei Gesetze:
% 1. Planetenbahnen sind Ellipsen mit der Sonne in einem Brennpunkt.
% 2. Der Fahrstrahl überstreicht in gleichen Zeiten gleiche Flächen.
% 3. T^2 \propto a^3.
% Wichtig wird das zweite Gesetz: es macht den Zeit-Position-
% Zusammenhang transzendent.

Die drei Kepler'schen Gesetze (1609, 1619) beschreiben die Planetenbewegung wie folgt:
\begin{enumerate}
\item Die Planetenbahnen sind Ellipsen, in deren einem Brennpunkt die Sonne steht.
\item Der Fahrstrahl von der Sonne zum Planeten überstreicht in gleichen Zeiten gleiche Flächen.
\item Die Quadrate der Umlaufzeiten verhalten sich wie die Kuben der grossen Halbachsen, $T^2 \propto a^3$.
\end{enumerate}
% [TODO: Bild einer Ellipse mit Brennpunkten und Halbachsen]

Während das erste und dritte Gesetz die Bahn rein \emph{geometrisch} und \emph{algebraisch} festlegen, sorgt das zweite Gesetz dafür, dass der Zusammenhang zwischen Zeit und Position auf der Bahn nicht algebraisch ausgedrückt werden kann.

\subsection{Drei Anomalien
\label{bessel:subsection:anomalien}}
% Inhalt:
% - Wahre Anomalie f (Winkel von der Sonne aus, also vom Brennpunkt).
% - Exzentrische Anomalie E (Winkel vom Mittelpunkt aus, auf dem
%   Hilfskreis um die Ellipse).
% - Mittlere Anomalie M = 2 pi t / T, linear in der Zeit.
% Mit einer Skizze einleuchtend.

In der Himmelsmechanik werden drei Winkel verwendet, um die Position eines Planeten auf seiner Bahn anzugeben:
\begin{itemize}
\item Die \emph{wahre Anomalie} $f$, gemessen vom Brennpunkt (also von der Sonne) aus.
\item Die \emph{exzentrische Anomalie} $E$, gemessen vom Mittelpunkt der Ellipse aus, auf dem die Ellipse umschreibenden Hilfskreis.
\item Die \emph{mittlere Anomalie} $M = 2\pi t / T$, die linear mit der Zeit wächst und der wahren Anomalie eines fiktiven gleichförmig umlaufenden Planeten entspricht.
\end{itemize}
% [TODO: Bild mit Ellipse, Hilfskreis, eingezeichneten Winkeln f, E, M]

\subsection{Die Kepler-Gleichung
\label{bessel:subsection:keplergleichung}}
% Inhalt:
% - Aus dem zweiten Kepler'schen Gesetz folgt durch geometrische
%   Überlegung: M = E - e sin(E).
% - Das ist transzendent in E -- nicht in geschlossener Form auflösbar.
% - Problem: gegeben M, finde E. Daraus dann f.

Aus dem Flächensatz folgt durch geometrische Betrachtung der Ellipse und ihres Hilfskreises die fundamentale Beziehung
\begin{equation}
\boxed{\;
M = E - e \sin E
\;}
\label{bessel:equation:kepler-gleichung}
\end{equation}
zwischen mittlerer und exzentrischer Anomalie, mit $e$ als numerischer Exzentrizität der Bahn.
% [TODO: Herleitung über Flächenvergleich]
Diese sogenannte \emph{Kepler-Gleichung} ist transzendent in $E$ und besitzt keine Lösung in geschlossener Form.
Das praktische Problem besteht darin, zu gegebenem $M$ (also gegebener Zeit $t$) die exzentrische Anomalie $E$ zu bestimmen.

\subsection{Fourier-Bessel-Entwicklung
\label{bessel:subsection:fourier-bessel}}
% Inhalt:
% - Die Funktion E(M) - M ist ungerade und 2-pi-periodisch in M,
%   bei festem e.
% - Daher besitzt sie eine Fourier-Sinusreihe:
%   E - M = \sum_{n=1}^{\infty} b_n(e) sin(n M).
% - Die Fourier-Koeffizienten berechnen sich aus dem Integral
%   b_n = (2/pi) \int_0^{pi} (E - M) sin(n M) dM,
%   das sich mit der Kepler-Gleichung umformen und in eine
%   Integraldarstellung von J_n(n e) überführen lässt.
% - Bessel zeigte 1824 die Identität b_n(e) = 2 J_n(n e) / n.
% - Endformel:
%   E = M + 2 \sum_{n=1}^{\infty} (J_n(n e)/n) sin(n M).
%
% Hier nicht alle Details ausrechnen, sondern den Weg skizzieren
% und den Anschluss an die Definition von J_n in Abschnitt 2 betonen.

Bei festem $e$ ist die Funktion $E - M$ als Funktion von $M$ ungerade und $2\pi$-periodisch.
Sie besitzt daher eine Fourier-Sinusreihe der Form
\begin{equation}
E - M
=
\sum_{n=1}^{\infty} b_n(e) \sin(nM).
\label{bessel:equation:fourier-ansatz}
\end{equation}
Die Fourier-Koeffizienten lassen sich durch ein Integral ausdrücken, das sich nach geschickter Umformung in die Integraldarstellung der Bessel-Funktion umwandelt.
% [TODO: Diese Rechnung detailliert ausführen oder skizzieren]
Bessel zeigte, dass
\begin{equation}
b_n(e) = \frac{2}{n}\, J_n(ne)
\label{bessel:equation:bessel-koeff}
\end{equation}
gilt.
Damit ergibt sich die berühmte Reihenentwicklung
\begin{equation}
\boxed{\;
E(M)
=
M + 2 \sum_{n=1}^{\infty} \frac{J_n(ne)}{n}\, \sin(nM).
\;}
\label{bessel:equation:kepler-loesung}
\end{equation}
Dies ist die historische Geburtsstelle der Bessel-Funktionen.

\subsection{Schlussbemerkung
\label{bessel:subsection:schluss}}
% Inhalt:
% Kurze Reflexion: zwei sehr unterschiedliche physikalische Probleme --
% das eine in der Wellenoptik, das andere in der Himmelsmechanik --
% führen auf dieselbe Funktionenklasse. Genau hier zeigt sich, was
% das Buch "Algebra und Analysis" als Leitthema verfolgt: algebraisch
% geprägte Strukturen (Kepler'sche Geometrie, Helmholtz-Operator in
% Koordinaten) treffen auf rein analytische Werkzeuge (Potenzreihen,
% Fourier-Reihen), und beides verschmilzt zu einer einzigen Familie
% spezieller Funktionen.

Mit der Beugung am Kreisspalt und der Kepler-Gleichung wurden zwei sehr unterschiedliche Phänomene betrachtet: das eine aus der Wellenoptik, das andere aus der Himmelsmechanik.
Beide führen auf dieselbe Familie spezieller Funktionen.
Hierin zeigt sich exemplarisch, was als Leitthema dieses Buches gilt: Die algebraisch geprägte Geometrie eines Problems -- sei es die Ellipsenbahn eines Planeten oder die kreisförmige Apertur einer Linse -- trifft auf analytische Werkzeuge wie den Potenzreihenansatz oder die Fourier-Reihe, und aus dieser Verbindung entsteht eine Funktionenklasse, die in beiden Welten gleichermassen unverzichtbar ist.

